% Chapter 5, section 3: Configuration for different use cases
% Write in plain LaTeX. Code of a guide step: \guide{<step>}; commands: \begin{shell}...\end{shell};
% a screenshot: \screenshot{file}{caption}; a gap in the draft: \placeholder{...}.
% See book/latex/queen.sty for the full list.
%
% Draft by Claude, 2026-10-23. Settings and defaults checked against server/src/config.rs and the
% site's operate and reference pages (webdoc/src/content/docs/operate/, reference/configuration.mdx).

\section{Configuration for different use cases}

Every setting is an environment variable, and most installations need very few of them. When a
node starts, it prints the settings it is really using, in lines that start with
\texttt{config:}. Read them after every change. A yes/no setting accepts \texttt{true} or
\texttt{false}, \texttt{1} or \texttt{0}, \texttt{yes} or \texttt{no}, \texttt{on} or
\texttt{off}; any other value stops the node at boot, so a typo cannot pass unnoticed.

\begin{pairs}[0.46]{You need}{What to set}
To try Queen, or to develop & Nothing \\
A small production & One node with a volume, and backups \\
To survive the loss of a machine & Three nodes; five to survive two \\
To keep strangers out & Broker JWT, or the built-in proxy \\
Many teams or customers on one cluster & The proxy, and raft groups \\
Payloads that must not sit on disk in clear & \texttt{QUEEN\_ENCRYPTION\_KEY} \\
Disk that does not fill up & Retention on the queues \\
Many short-lived partitions & \texttt{PARTITION\_CLEANUP\_SECONDS} and two more \\
Kafka clients & \texttt{QUEEN\_KAFKA\_EMBEDDED} \\
Thousands of queues & \texttt{QUEEN\_QLOG\_SHARDS}, before the first start \\
To survive the loss of a whole site & A standby cluster \\
\end{pairs}

\subsection{To try Queen, or to develop}

One command, no settings:

\begin{shell}
docker run -p 6632:6632 ghcr.io/queen-mq/queen:2.6.0
\end{shell}

Without a volume, the data lives inside the container and is gone when the container is removed.
That is usually what you want on a laptop.

\subsection{A small production}

One node with its data on a volume:

\begin{shell}
docker run -d --name queen --restart unless-stopped --stop-timeout 60 \
  -p 6632:6632 -v /var/lib/queen:/var/lib/queen/raft \
  ghcr.io/queen-mq/queen:2.6.0
\end{shell}

A single node writes every message to disk (fsync) before it answers. A crash, a kill -9 or a
power cut loses nothing it acknowledged: it starts again and reads its log. What it cannot survive
is losing its disk, so take backups (see Recovery). An upgrade stops it for a few seconds, and its
ephemeral queues, which live in memory, are empty after every restart. Port 6632 stays on a private
network, as on a cluster (see To keep strangers out).

\why{many workloads do not need three machines. One node is the simplest thing that keeps every
acknowledged message. Choose it knowing that a lost disk means restoring the last backup.}

\subsection{To survive the loss of a machine}

Run three nodes, or five (see Run a cluster with Docker). A majority of three is two, so three
nodes survive one failure. Five survive two, but every write then waits for three disks, and the
leader feeds four followers. Four nodes survive only one failure, like three, while every write
waits for three of them: an even number only costs. Put each node on its own machine, and in its
own zone if you have zones.

\why{a majority is what lets the cluster decide alone, without a person, which node leads. More
nodes make a cluster harder to kill, not faster: one leader orders every write, so adding nodes
adds no write throughput.}

\subsection{To keep strangers out}

Out of the box a node checks nobody: anyone who reaches port 6632 can read, write, delete queues
and change the cluster's members. Choose one of two front doors.

\queenfig{front-doors}{What listens where: only the proxy's port faces the internet.}

\begin{options}{Setting}{Default}{What it does}
JWT\_ENABLED & false & The broker checks a token on every request. For services that already get tokens from an identity provider. \\
JWT\_ALGORITHM & HS256 & HS256 checks with \texttt{JWT\_SECRET}; RS256 and EdDSA with \texttt{JWT\_PUBLIC\_KEY} or \texttt{JWT\_JWKS\_URL}. \\
QUEEN\_PROXY\_EMBEDDED & false & The built-in proxy: API keys, people who sign in to the dashboard, tenants, quotas. \\
QUEEN\_PROXY\_PORT & unset & The proxy's own port. Expose this one, and keep 6632 inside. \\
QUEEN\_TLS\_CERT, \_KEY & unset & TLS on the proxy's port. The broker port has no TLS: use your load balancer. \\
\end{options}

A token's roles decide what it may do: \texttt{read-only}, \texttt{write-only} (push, for producers
that must not read), \texttt{read-write} and \texttt{admin} (the membership API and deletes).
Whatever you choose, keep the raft port (7400) on a private network, with \texttt{QUEEN\_RAFT\_TOKEN}
set on every node.

\why{a broker holds every payload and every key of your system. The default is open so that the
first run is one command, and it is the first thing to change before anything else can reach the
node.}

\subsection{Many tenants on one cluster}

One cluster can serve many teams or customers. Each one is a tenant: its queues, KV and timers are
invisible to the others. The proxy built into the broker adds API keys and people in front of it.
It keeps its own data (tenants, users, key hashes, plans, usage) in the broker's replicated KV, so
there is no database to run.

\begin{shell}
SESSION_SECRET=$(openssl rand -hex 32)    # signs console sessions
CP_TOKEN=$(openssl rand -hex 32)          # opens the control plane, /api/cp/*
QUEEN_KEY="qk_live_$(openssl rand -base64 32 | tr '+/' '-_' | tr -d '=')"

docker run -d --name queen -p 6711:6711 -p 127.0.0.1:6632:6632 \
  -v /var/lib/queen:/var/lib/queen/raft \
  -e QUEEN_PROXY_EMBEDDED=true \
  -e QUEEN_PROXY_PORT=6711 \
  -e QUEEN_PROXY_JWT_SECRET="$SESSION_SECRET" \
  -e QUEEN_PROXY_CP_TOKEN="$CP_TOKEN" \
  -e QUEEN_PROXY_SPOOL_DIR=/var/lib/queen/raft/proxy-spool \
  -e QUEEN_PROXY_BOOTSTRAP_TENANT=acme \
  -e QUEEN_PROXY_BOOTSTRAP_API_KEY="$QUEEN_KEY" \
  -e QUEEN_PROXY_DEFAULT_CLUSTER=acme \
  ghcr.io/queen-mq/queen:2.6.0
\end{shell}

The first start creates the tenant \texttt{acme} and its key. Clients now use port 6711, with the
key (\texttt{bearerToken} in the JavaScript client); without it they get \texttt{401
unauthorized}. Port 6632 is the broker itself, with no keys: the nodes talk to each other on it,
so it stays inside your network.

Add a tenant with one call to the control plane. You make the key yourself and send only its
hash, so the proxy never sees the key:

\begin{shell}
KEY="qk_live_$(openssl rand -base64 32 | tr '+/' '-_' | tr -d '=')"
HASH=$(printf %s "$KEY" | sha256sum | cut -d' ' -f1)
curl -s -X POST localhost:6711/api/cp/provision -H "x-queen-cp-token: $CP_TOKEN" \
  -H 'content-type: application/json' -d '{
    "tenant_slug": "globex", "cluster_slug": "globex", "plan": "pro",
    "user_id": "308d3f3d-2306-41de-95f1-e32b12ec7470", "email": "ops@globex.example",
    "role": "admin", "key_name": "default", "key_hash": "'"$HASH"'",
    "scopes": ["produce", "consume", "read"]}'
{"broker_tenant_uuid":"91a5d7bb-879c-4ae2-95e2-295fa2704098","created":{"cluster":true,...}}
\end{shell}

The proxy finds the tenant from the first part of the host name: \texttt{globex.\allowbreak queen.\allowbreak example.com}
is \texttt{globex}. A key works only there: on another tenant's host it gets \texttt{403 key/cluster
mismatch}. Each plan (\texttt{free}, \texttt{dev}, \texttt{pro}, \texttt{dedicated-s}) limits
requests a second, queues and stored bytes. The proxy starts by only logging what it would block;
\texttt{QUEEN\_PROXY\_ENFORCE=true} makes it block.

With one raft group, all tenants share one leader. \texttt{QUEEN\_RAFT\_GROUPS=3} runs three groups
in each node, each with its own leader, log and directory (\texttt{groups/g1/}, \ldots). A tenant
lives in one group, chosen by a hash of its id. To choose yourself, pin it by the
\texttt{broker\_tenant\_uuid} of the answer above: \texttt{QUEEN\_TENANT\_GROUPS=<uuid>=1} puts it in
group 1. Pin a big tenant before it writes, because a tenant cannot move while it has data. With as many groups as nodes, every node leads one group.

\queenfig{tenants-groups}{Tenants through the proxy, each in one raft group.}

\why{tenants share machines, not data. One leader per tenant keeps each tenant's transaction a
single log entry, and more groups spread the tenants' writes over more leaders. Raft groups were
not part of the Jepsen tests, which ran one group.}

Three things to know: every node needs the same secrets; the membership routes are closed on the
proxy's port, so call them on 6632; and keep \texttt{QUEEN\_PROXY\_SPOOL\_DIR} on the data volume,
where usage waits when it cannot be written.

\subsection{Payloads that must not sit on disk in clear}

Give every node the same key, 64 hex characters, then turn encryption on for the queues that need
it (\texttt{encryptionEnabled}, see Create and configure queues):

\begin{shell}
QUEEN_ENCRYPTION_KEY="$(openssl rand -hex 32)"      # the same value on every node
\end{shell}

Payloads are encrypted with AES-256-GCM. Queue and partition names, KV keys and KV values are not.
Check that the boot log says \texttt{encryption service initialized (AES-256-GCM)}. Without a key, the
node stores those queues in clear and says so only in its \texttt{config:} lines; a key of the wrong
length logs one warning and does the same.

\why{some payloads hold personal data, and a disk, or a copy of it, can end up in the wrong hands.
The price is the key itself: lose it and the payloads are gone; change it and consumers receive the
encrypted envelope as their data, with no error. Keep it safe, apart from the backups.}

\subsection{Disk that does not fill up}

Queen keeps every message until retention removes it, even after every group has acked it (see
Create and configure queues). Turn retention on for the queues that grow. As a last guard, a node
whose disk is 85\% full (\texttt{QUEEN\_RAFT\_DISK\_HIGH\_PCT}) refuses the writes that would grow
it, with \texttt{507 storage\_full}, and accepts them again below 80\%
(\texttt{QUEEN\_RAFT\_DISK\_LOW\_PCT}). Reads, acks and deletes keep working.

\why{a broker that keeps history lets new consumer groups and replays read old messages. Retention
is where you decide how long that history is worth its disk. The 85\% gate leaves room to grow the
volume before the node stops.}

\subsection{Many short-lived partitions}

Some workloads open a partition per trace, per session or per order, and never write to it again.
By default an empty partition stays for 30 days, and each one costs memory and time in every
retention round. To delete them soon after their last message, set three things on the broker:

\begin{shell}
docker run -d --name queen -p 6632:6632 -v /var/lib/queen:/var/lib/queen/raft \
  -e PARTITION_CLEANUP_SECONDS=60 \
  -e QUEEN_RAFT_TXN_WINDOW_MIN_S=60 \
  -e QUEEN_RAFT_RETENTION_SCAN_PER_S=40000 \
  ghcr.io/queen-mq/queen:2.6.0
\end{shell}

and give the queue retention and a short dedup window:

\begin{shell}
curl -s -X POST http://localhost:6632/api/v1/configure -H 'content-type: application/json' \
  -d '{"queue":"traces","options":{"retentionEnabled":true,
       "completedRetentionSeconds":60,"dedupWindowSeconds":60}}'
\end{shell}

Retention removes a message once it is a minute old and every group has acked it, and the empty
partition is deleted after a minute without a write. We pushed 50 messages to 50 partitions and
consumed them: all 50 partitions were gone within 76 seconds.

\queenfig{short-lived-partitions}{A short-lived partition: three settings, and the longest wins.}

\why{a partition holds the dedup records of its messages. Deleted too early, a producer that
retries would store a message twice. So a partition is never deleted inside its queue's dedup
window, nor before \texttt{QUEEN\_RAFT\_TXN\_WINDOW\_MIN\_S} (15 minutes by default): a short cleanup
needs both to be short too.}

On 2.6.0 the third setting matters when partitions come and go fast: we saw the problem from 200
new partitions a second. By default the retention scan paces each round from the size of the previous one. When new
partitions arrive faster than that, a round never ends, and nothing is deleted: in a test with
1,000 new partitions a second, 594,000 of 595,000 partitions were still there after ten minutes.
\texttt{40000} turns that pacing off. With it, 1,000 and 2,000 new partitions a second stayed flat,
at about 70,000 and 150,000 partitions.

\subsection{Kafka clients}

\texttt{QUEEN\_KAFKA\_EMBEDDED=true} opens a Kafka port (9092) on every node, so existing Kafka
producers and consumers work without code changes. On a cluster, each node also needs
\texttt{QUEEN\_KAFKA\_ADVERTISED\_ADDR}, the address at which clients reach that node's Kafka
port. The Kafka port has no TLS in 2.6, and passwords (\texttt{QUEEN\_KAFKA\_SASL=plain}) cross the
network in clear: keep it on a private network. How to use with Kafka clients, in chapter 4, has
the details.

\subsection{Thousands of queues}

By default each queue has its own log file on disk, and a write can touch many files.
\texttt{QUEEN\_QLOG\_SHARDS=K} puts all queues into K shared logs instead, so a write syncs at most
K files. It is read only when the data directory is created, so set it before the first start.

\subsection{To survive the loss of a whole site}

Three nodes in one zone, or in one Kubernetes cluster, all go down together when that goes. A
standby is a second cluster, somewhere else, that reads everything the first one commits and stays
a moment behind it. When the first is lost, you promote the standby and point the clients at it.

\why{replicas protect against a lost machine, backups against a mistake. A standby protects
against losing everything at once, and loses only the last moment before the failure. It does not
protect against mistakes: a deleted queue is replayed on the standby a moment later. Keep the
backups too.}

On the site: \href{https://queenmq.com/manage/configuration/}{queenmq.com/manage/configuration}. More there: \href{https://queenmq.com/reference/configuration/}{every setting}, \href{https://queenmq.com/manage/security/}{security}, \href{https://queenmq.com/manage/tenants/}{tenants and API keys}, \href{https://queenmq.com/manage/standby/}{standby cluster}.
